%
% Slide template for the FreeBSD Developer Summit.
% Take it or leave it :-)
%
% It requires LaTeX and LaTeX Beamer [1] to compile.
% pdfLaTeX is recommended for compilation as it produces
% PDF file immediately.
%
% $ pdflatex some.latex
%
%
% It is also recommended to convert images to PDF
% by using ImageMagick ("convert") before including them.
%
% [1] http://latex-beamer.sourceforge.net
%

\documentclass{beamer}

\usepackage{url}
\usepackage[utf8]{inputenc}
\usepackage[russian]{babel}
\usepackage{verbatim}
\usepackage{graphicx}
\usepackage{listings}
\usepackage{mathtools}
\usepackage{color}

\mode<presentation>
{
  \definecolor{beamer@gker}{rgb}{0.8,0.0,0.0}
  \setbeamercolor*{structure}{fg=beamer@gker}
  \logo{\includegraphics[scale=0.5]{logo.pdf}}
}

\setbeamertemplate{footline}[text line]{%
  \parbox{\linewidth}{\vspace*{-8pt}
  \url{http://highsecure.ru/pkg}
  \hfill
  \insertshortauthor
  \hfill
  \insertframenumber\ of \inserttotalframenumber}}
\setbeamertemplate{navigation symbols}{}

\definecolor{light-gray}{gray}{0.60}

\title{Pkg - система управления пакетами FreeBSD}
\author{Всеволод Стахов \\ \url{vsevolod@FreeBSD.org}}
\institute{\includegraphics[scale=0.5]{logo.pdf}}
\date{ruBSD conference 14 Декабря 2013}

\begin{document}

\begin{frame}[plain]
  \titlepage
\end{frame}

\begin{frame}
\frametitle{Порты и пакеты}
Порты - основа для создания пакетов.
\begin{itemize}
  \item Устоявшаяся система
  \item Понятные и четкие правила
  \item Простые в создании и управлении (не всегда)
  \item Множество настроек
\end{itemize}
\end{frame}

\begin{frame}[fragile]
\frametitle{Порты без pkg}
\begin{figure}[h!]
  \centering
  \includegraphics[width=0.95\textwidth]{q1.eps}
\end{figure}
\end{frame}

\begin{frame}
\frametitle{Недостатки предыдущей схемы}
\begin{itemize}
  \item Make не может полноценно разрешать зависимости и конфликты
  \item Сложно обновлять и держать в актуальном состоянии
  \item Проблемы при миграции между версиями ОС
  \item Процесс компиляции из исходников занимает время
\end{itemize}
\end{frame}

\begin{frame}
\frametitle{Планируемая архитектура взаимодействия pkg и портов}
\begin{figure}[h!]
  \centering
  \includegraphics[width=0.99\textwidth]{q2.eps}
\end{figure}
\end{frame}

\begin{frame}
\frametitle{Порты и пакеты}
\begin{itemize}
  \item Порты используются для построения пакетов
  \item Зависимости и конфликты обрабатываются pkg, а не make
  \item Стабильная ветка портов соответствует стабильной ветке пакетов
  \item Пользователям рекоммендуется ставить ПО из пакетов
  \item Не запрещая установку из портов
\end{itemize} 
\end{frame}

\begin{frame}
\frametitle{Создание репозитория пакетов}
\begin{figure}[h!]
  \centering
  \includegraphics[width=0.95\textwidth]{q3.eps}
\end{figure}
\end{frame}

\begin{frame}
\frametitle{Установка пакетов в систему}
\begin{figure}[h!]
  \centering
  \includegraphics[width=0.95\textwidth]{q4.eps}
\end{figure}
\end{frame}

\begin{frame}
\frametitle{Проблемы в pkg}
\begin{itemize}
  \item Поддержка устаревших портов (например, без stage directory)
  \item Крайне ограниченная система зависимостей
  \item Примитивный алгоритм разрешения зависимостей (solver)
\end{itemize}
\end{frame}

\begin{frame}
\frametitle{Проблемы с алгоритмом solver'а}

\begin{itemize}
\item Нет поддержки конфликтов
\item Не работает с альтернативными пакетами
\item Может выполнять только одну задачу: установить, удалить или обновить
некоторые пакеты
\end{itemize}

\end{frame}


\begin{frame}
\frametitle{Существующие системы управления пакетами}

Другие системы управления пакетами используют различные подходы для решения
данной задачи:
\begin{itemize}
  \item \includegraphics[width=22pt]{suse.pdf} \hspace{5pt} \texttt{Zypper/SUSE}
  - построена на базе libsolv
  \item \includegraphics[width=22pt]{rpm.pdf} \hspace{5pt} \texttt{Yum/RedHat} -
  планирует переход на libsolv
  \item \includegraphics[width=15pt]{debian.pdf} \hspace{12pt}
  \texttt{Apt/Debian} - базируется на собственной логике разрешения зависимостей
  и конфликтов
  \item \includegraphics[width=15pt]{arch.pdf} \hspace{12pt}
  \texttt{Pacman/Archlinux} - использует собственный наивный алгоритм
\end{itemize}

\end{frame}

\begin{frame}[fragile]
\frametitle{Независимые solver'ы}
Для взаимодействия с внешними solver'ами, предполагается использовать формат
CUDF, разработанный в рамках исследовательского проекта Mancoosi
\url{http://mancoosi.org}:
\bigskip
{\small
	\begin{verbatim}
	package: devel/libblah
	version: 1
	depends: x11/libfoo

	package: security/blah
	version: 2
	depends: devel/libblah
	conflicts: security/blah-devel
	
	\end{verbatim}
}
\end{frame}

\begin{frame}
\frametitle{Но нам нужен собственный solver!}

Варианты:
\begin{itemize}
  \item Реализовать собственный алгоритм разрешения зависимостей?
  \pause
  \item Использовать существующие решения?
  \pause
  \item Использовать известный алгоритм?
  \pause
\end{itemize}
\bigskip
{\large SAT solver для управления пакетами}
\bigskip
\[
\overbrace{\underbrace{(x_1 \| \neg x_2 \|
x_3)}_{\text{\color{light-gray}Условие}} \& (x_3 \| \neg x_1) \&
(x_2)}^{\text{\color{light-gray}SAT выражение}}
\]
\end{frame}

\begin{frame}
\frametitle{Составление выражения SAT}
\begin{itemize}
  \item Каждой версии каждого пакета назначается независимая переменная: 
  A $\to a_1$, B $\to b_1$
  \item Пользовательский запрос преобразуется в набор унарных условий:
  \begin{itemize}
    \item Установить или обновить пакет A $\to (a_1)$
    \item Удалить пакет B $\to (\neg b_1)$
  \end{itemize}
  \item Преобразовать правила зависимостей и конфликтов в условия
\end{itemize}

\end{frame}

\begin{frame}
\frametitle{Преобразование конфликтов и зависимостей}
\begin{itemize}
  \item Если пакет A зависит от пакета B (версий $B_1$ и $B_2$), тогда либо
  пакет А не установлен, либо установлена одна из версий пакета B:
  \bigskip
\[
(\neg A \| B_1 \| B_2)
\]
\pause
  \item Если разные версии пакета B ($B_1$, $B_2$ and $B_3$) конфликтуют между
  собой, тогда только одна из версий пакета B может быть установлена:
  \bigskip
\[
\underbrace{(\neg B_1 \| \neg B_2) \& (\neg B_1 \| \neg B_3) \& (\neg B_2 \|
\neg B_3)}_{\text{\color{light-gray}Цепочка конфликтов}}
\]
\end{itemize}
\end{frame}

\begin{frame}
\frametitle{Алгоритм решения задачи SAT}

Способы упростить решение задачи SAT.
\begin{itemize}
  \item Очевидные назначения
  \item Назначения юнитов - назначение переменных в условиях, содержащих только
  одну неназначенную переменную и не являющихся истинными
  \item Обучение на базе конфликтов: если в ходе назначения мы обнаружили
  конфликт, то двигаемся по дереву вверх, до тех пор, пока все выражение
  является ложным, после чего инвертируем назначение, вызвавшее конфликт
  \item Предположения, специфичные для задачи управления пакетами
\end{itemize}
\end{frame}

\begin{frame}
\frametitle{Назначения переменных}
\begin{itemize}
  \item Очевидные назначения - переменные, непосредственно входящие в запрос
  пользователя (установка или удаление)
  \bigskip
  \[
  (\neg A \| B) \& \underbrace{(A)}_{\textit{\color{light-gray}true}} \&
  \underbrace{(\neg C)}_{\textit{\color{light-gray}false}} \& (\neg A \| \neg D)
  \]
  \pause
  \item Назначения юнитов - простые зависимости и конфликты
  \bigskip
  \[
  \overbracket{\underbrace{(\neg A \| B)}_{B \rightarrow
  \textit{\color{light-gray}true}}}^{\text{\color{light-gray}Зависимость}} \&
  \overbrace{(A)}^{\textit{\color{light-gray}true}} \& \overbrace{(\neg
  C)}^{\textit{\color{light-gray}false}} \&
  \overbracket{\underbrace{(\neg A \| \neg D)}_{D \rightarrow
  \textit{\color{light-gray}false}}}^{\text{\color{light-gray}Конфликт}}
  \]
  \
\end{itemize}
\end{frame}

\begin{frame}
\frametitle{Обучение на основание конфликтов}
Для обработки альтренативных решений необходимо перебирать все возможные
варианты назначений:
\begin{enumerate}
  \item полный перебор в глубину
  \item возврат, если был обнаружен конфликт
  \item определить, какое назначение вызвало конфликт
  \item инвертировать такое назначение и продолжить поиск, отсекая конфликтную
  ветвь
\end{enumerate}
\end{frame}

\begin{frame}
\frametitle{Особенности задачи управления пакетами}
Решение задачи SAT в чистом виде может быть сильно ускорено, если принимать во
внимание специфику исходной задачи:
\begin{itemize}
  \item не пытаться удалить уже установленные пакеты (если не обнаружено
  конфликтов)
  \item не устанавливать пакеты, если они явно не нужны
  \item предпочитать установку более приоритетных пакетов, учитывая приоритет
  репозиториев
\end{itemize}
Эти предположения также улучшают начальное назначение, отсекая заведомо
неверные ветви дерева.
\end{frame}

\begin{frame}
\frametitle{Множество пакетов}
Мы представляем все пакеты, вовлеченные в запрос пользователя (все прямые и
обратные зависимости, а также конфликты) в виде множества пакетов:
\begin{figure}[h!]
  \centering
  \includegraphics[height=0.7\textheight]{q5.eps}
\end{figure}
\end{frame}

\begin{frame}
\frametitle{Задача управления пакетами}
\begin{itemize}
  \item Запрос пользователя превращается в набор задач по утсановке или удалению
  пакетов
  \item Все явные и неявные конфликты проверяются во время создания репозитория
  \item Создается множество пакетов на основании всех зависимостей и конфликтов
  пакетов, входящих в запрос
  \item Множество пакетов строится на базе имени и уникальной версии, которая
  определяется всеми значащими полями пакета (версия, ревизия, контрольные
  суммы)
\end{itemize}
\end{frame}

\begin{frame}
\frametitle{Применение алгоритмов solver'ов в pkg}
\begin{itemize}
  \item Pkg может использовать множество пакетов для работы с внешним solver'ом:
  \begin{itemize}
    \item формализация версий
    \item генерация запроса
    \item обработка решения
  \end{itemize}
  \item Также можно использовать встроенный SAT solver:
  \begin{itemize}
    \item преобразовать множество пакетов в булево выражение
    \item присоединить запрос
    \item ???
    \item PROFIT
  \end{itemize} 
\end{itemize}
\end{frame}

\begin{frame}
\frametitle{Возможности}

\begin{itemize}
  \item Использование pkg для управления портами
  \item Лучшая поддержка нескольких репозиториев
  \item Проверка алгоритмов управления пакетами (через CUDF формат)
  \item Продвинутый формат зависимостей и конфликтов
  \item Альтернативные пакеты
\end{itemize}
\end{frame}

\begin{frame}
\frametitle{Новый формат зависимостей}
\large{\[libblah >= 1.0 +option_1, +option_2 \| libfoo != 1.1\]}
\begin{itemize}
  \item Зависимости от реальных и виртуальных пакетов
  \item Более гибкий вариант определения версий
  \item Использование альтернатив
\end{itemize}
\begin{figure}[h!]
  \centering
  \includegraphics[height=0.3\textheight]{q6.eps}
\end{figure}
\end{frame}

\begin{frame}
\frametitle{Альтернативные пакеты}
\begin{itemize}
  \item Используются для определения некоторой общей функциональности
  (например, браузер)
  \item Могут использоваться в зависимостях (т.н. виртуальные зависимости)
\end{itemize}
\begin{figure}[h!]
  \centering
  \includegraphics[height=0.4\textheight]{q7.eps}
\end{figure}
\end{frame}

\begin{frame}
\begin{center}
\includegraphics{logo.pdf}
{\Large Спасибо за внимание!} \\
\emph{Вопросы?} \\[4pt]
\url{vsevolod@FreeBSD.org}
\end{center}
\end{frame}

\end{document}
